\documentclass[12pt]{ximera}
\title{Homework 6: Independence, Bayes, and the Multiplication Principle}
\begin{document}
\begin{abstract}
Sections 7.5, 7.6, and 8.1: recognizing independent events, total probability and
Bayes' Theorem from a voting example, and counting arrangements with the
multiplication principle.
\end{abstract}
\maketitle

\section*{Sections 7.5--7.6 and 8.1}

\begin{problem}
Identify whether each of the following pairs of events is independent or not.

\begin{prompt}
\textbf{(a)} Event $E$: your laundry machine is broken. \\
Event $F$: you have worn the same shirt multiple times in the last week.

\begin{multipleChoice}
\choice{Independent}
\choice[correct]{Not independent}
\end{multipleChoice}

\begin{feedback}
Not independent. A broken machine makes re-wearing a shirt more likely, so
$P(F\mid E)>P(F)$.
\end{feedback}
\end{prompt}

\begin{prompt}
\textbf{(b)} Drawing cards from a deck \emph{with} replacement. \\
Event $E$: you draw a red card on the first draw. \\
Event $F$: you draw a red card on the second draw.

\begin{multipleChoice}
\choice[correct]{Independent}
\choice{Not independent}
\end{multipleChoice}

\begin{feedback}
Independent. The first card is returned, so the second draw faces the full deck either
way: $P(F\mid E)=\frac{26}{52}=P(F)$.
\end{feedback}
\end{prompt}

\begin{prompt}
\textbf{(c)} Drawing cards from a deck \emph{without} replacement. \\
Event $E$: you draw a red card on the first draw. \\
Event $F$: you draw a red card on the second draw.

\begin{multipleChoice}
\choice{Independent}
\choice[correct]{Not independent}
\end{multipleChoice}

\begin{feedback}
Not independent. After a red card is removed, $P(F\mid E)=\frac{25}{51}$, which differs
from $P(F)=\frac12$.
\end{feedback}
\end{prompt}

\begin{prompt}
\textbf{(d)} Event $E$: the high today is above $60^\circ$ Fahrenheit. \\
Event $F$: you are dressed in a sweater.

\begin{multipleChoice}
\choice{Independent}
\choice[correct]{Not independent}
\end{multipleChoice}

\begin{feedback}
Not independent. A warm day makes a sweater less likely, so $P(F\mid E)<P(F)$.
\end{feedback}
\end{prompt}
\end{problem}

\begin{problem}
Use a tree diagram and/or Bayes' Theorem to solve the following.

Suppose in an election, only $60\%$ of eligible voters actually voted. Among those who
voted, $35\%$ were registered Republican, $30\%$ Democrat, $25\%$ Independent, and $10\%$
other; among those who didn't vote, $20\%$ were registered Republican, $25\%$ Democrat,
$30\%$ Independent, and $25\%$ other.

\begin{prompt}
\textbf{(a)} What is the total probability that an eligible voter was registered
Independent?

$P(I) = \answer[tolerance=0.001]{0.27}$

\begin{hint}
Two branches of the tree end at ``Independent'': voted-and-Independent, and
didn't-vote-and-Independent. Add them.
\end{hint}

\begin{feedback}
\[
P(I)=P(V)P(I\mid V)+P\!\left(\olsi{V}\right)P\!\left(I\mid\olsi{V}\right)
=(0.60)(0.25)+(0.40)(0.30)=0.15+0.12=0.27.
\]
\end{feedback}
\end{prompt}

\begin{prompt}
\textbf{(b)} What is the probability that someone registered Independent voted in the
election --- that is, $P(V\mid I)$? Enter a fraction in lowest terms.

$\answer{5/9}$

\begin{feedback}
\[
P(V\mid I)=\frac{P(V\cap I)}{P(I)}=\frac{0.15}{0.27}=\frac{15}{27}=\frac59\approx 0.556.
\]
\end{feedback}
\end{prompt}
\end{problem}

\begin{problem}
Use the multiplication principle to count the following.

A family has ten siblings, and it is hard to remember their birth order perfectly.

\begin{prompt}
\textbf{(a)} Guessing at random, how many orderings of the ten siblings are there,
oldest to youngest?

$\answer{3628800}$

\begin{feedback}
There are $10$ choices for the oldest, then $9$ for the next, and so on:
\[ 10\cdot 9\cdot 8\cdots 1 = 10! = 3{,}628{,}800. \]
\end{feedback}
\end{prompt}

\begin{prompt}
\textbf{(b)} Suppose you remember the positions of four of the siblings: Ann is oldest,
Bill is second, Carol is second-youngest, and Dave is youngest. How many orderings
have these four in their correct positions?

$\answer{720}$

\begin{feedback}
Four positions are settled, leaving the other $6$ siblings to fill the $6$ middle
positions in any order: $6! = 720$.
\end{feedback}
\end{prompt}

\begin{prompt}
\textbf{(c)} Suppose instead you remember only the pattern, oldest to youngest:
\[
\text{girl, boy, boy, boy, girl, boy, girl, boy, girl, boy}.
\]
How many orderings of the ten siblings fit this pattern, given there are $6$ boys and
$4$ girls?

$\answer{17280}$

\begin{hint}
The pattern fixes which positions hold girls and which hold boys. Arrange the girls
among their positions, and independently the boys among theirs.
\end{hint}

\begin{feedback}
The $4$ girls fill the four girl positions in $4!=24$ ways, and independently the $6$
boys fill the six boy positions in $6!=720$ ways:
\[ 4!\cdot 6! = 24\cdot 720 = 17{,}280. \]
\end{feedback}
\end{prompt}
\end{problem}

\begin{problem}
Suppose you want to plan outfits for a $7$-day trip. You've made a checklist of the
decisions you need to make:
\begin{itemize}
\item You own $20$ tops, and need to select and order $7$ for the trip.
\item You own $12$ bottoms, and similarly need to select and order $7$ for the trip.
\item You own $15$ pairs of socks, and need to select and order $7$ pairs for the trip.
\end{itemize}

\begin{prompt}
Write a formula for the number of possible outfit arrangements for the week.

\begin{freeResponse}
\end{freeResponse}

\begin{feedback}
Each checklist item selects \emph{and orders} $7$ items, so each is a permutation, and
the three are independent decisions:
\[
P(20,7)\cdot P(12,7)\cdot P(15,7) = \frac{20!}{13!}\cdot\frac{12!}{5!}\cdot\frac{15!}{8!}.
\]
\end{feedback}
\end{prompt}

\begin{prompt}
Enter the value of each factor:

$P(20,7)=\answer{390700800}$ \qquad $P(12,7)=\answer{3991680}$ \qquad
$P(15,7)=\answer{32432400}$

\begin{feedback}
\[
P(20,7)=20\cdot 19\cdots 14=390{,}700{,}800,\quad
P(12,7)=12\cdot 11\cdots 6=3{,}991{,}680,\quad
P(15,7)=15\cdot 14\cdots 9=32{,}432{,}400.
\]
Their product is about $5.06\times 10^{22}$ possible weeks.
\end{feedback}
\end{prompt}
\end{problem}

\end{document}
